% Options for packages loaded elsewhere
\PassOptionsToPackage{unicode}{hyperref}
\PassOptionsToPackage{hyphens}{url}
\PassOptionsToPackage{dvipsnames,svgnames,x11names}{xcolor}
%
\documentclass[
  10pt,
]{article}
\usepackage{amsmath,amssymb}
\usepackage{iftex}
\ifPDFTeX
  \usepackage[T1]{fontenc}
  \usepackage[utf8]{inputenc}
  \usepackage{textcomp} % provide euro and other symbols
\else % if luatex or xetex
  \usepackage{unicode-math} % this also loads fontspec
  \defaultfontfeatures{Scale=MatchLowercase}
  \defaultfontfeatures[\rmfamily]{Ligatures=TeX,Scale=1}
\fi
\usepackage[]{fontspec}
\ifPDFTeX\else
  % xetex/luatex font selection
\fi
% Use upquote if available, for straight quotes in verbatim environments
\IfFileExists{upquote.sty}{\usepackage{upquote}}{}
\IfFileExists{microtype.sty}{% use microtype if available
  \usepackage[]{microtype}
  \UseMicrotypeSet[protrusion]{basicmath} % disable protrusion for tt fonts
}{}
\makeatletter
\@ifundefined{KOMAClassName}{% if non-KOMA class
  \IfFileExists{parskip.sty}{%
    \usepackage{parskip}
  }{% else
    \setlength{\parindent}{0pt}
    \setlength{\parskip}{6pt plus 2pt minus 1pt}}
}{% if KOMA class
  \KOMAoptions{parskip=half}}
\makeatother
\usepackage{xcolor}
\usepackage{color}
\usepackage{fancyvrb}
\newcommand{\VerbBar}{|}
\newcommand{\VERB}{\Verb[commandchars=\\\{\}]}
\DefineVerbatimEnvironment{Highlighting}{Verbatim}{commandchars=\\\{\}}
% Add ',fontsize=\small' for more characters per line
\usepackage{framed}
\definecolor{shadecolor}{RGB}{248,248,248}
\newenvironment{Shaded}{\begin{snugshade}}{\end{snugshade}}
\newcommand{\AlertTok}[1]{\textcolor[rgb]{0.94,0.16,0.16}{#1}}
\newcommand{\AnnotationTok}[1]{\textcolor[rgb]{0.56,0.35,0.01}{\textbf{\textit{#1}}}}
\newcommand{\AttributeTok}[1]{\textcolor[rgb]{0.13,0.29,0.53}{#1}}
\newcommand{\BaseNTok}[1]{\textcolor[rgb]{0.00,0.00,0.81}{#1}}
\newcommand{\BuiltInTok}[1]{#1}
\newcommand{\CharTok}[1]{\textcolor[rgb]{0.31,0.60,0.02}{#1}}
\newcommand{\CommentTok}[1]{\textcolor[rgb]{0.56,0.35,0.01}{\textit{#1}}}
\newcommand{\CommentVarTok}[1]{\textcolor[rgb]{0.56,0.35,0.01}{\textbf{\textit{#1}}}}
\newcommand{\ConstantTok}[1]{\textcolor[rgb]{0.56,0.35,0.01}{#1}}
\newcommand{\ControlFlowTok}[1]{\textcolor[rgb]{0.13,0.29,0.53}{\textbf{#1}}}
\newcommand{\DataTypeTok}[1]{\textcolor[rgb]{0.13,0.29,0.53}{#1}}
\newcommand{\DecValTok}[1]{\textcolor[rgb]{0.00,0.00,0.81}{#1}}
\newcommand{\DocumentationTok}[1]{\textcolor[rgb]{0.56,0.35,0.01}{\textbf{\textit{#1}}}}
\newcommand{\ErrorTok}[1]{\textcolor[rgb]{0.64,0.00,0.00}{\textbf{#1}}}
\newcommand{\ExtensionTok}[1]{#1}
\newcommand{\FloatTok}[1]{\textcolor[rgb]{0.00,0.00,0.81}{#1}}
\newcommand{\FunctionTok}[1]{\textcolor[rgb]{0.13,0.29,0.53}{\textbf{#1}}}
\newcommand{\ImportTok}[1]{#1}
\newcommand{\InformationTok}[1]{\textcolor[rgb]{0.56,0.35,0.01}{\textbf{\textit{#1}}}}
\newcommand{\KeywordTok}[1]{\textcolor[rgb]{0.13,0.29,0.53}{\textbf{#1}}}
\newcommand{\NormalTok}[1]{#1}
\newcommand{\OperatorTok}[1]{\textcolor[rgb]{0.81,0.36,0.00}{\textbf{#1}}}
\newcommand{\OtherTok}[1]{\textcolor[rgb]{0.56,0.35,0.01}{#1}}
\newcommand{\PreprocessorTok}[1]{\textcolor[rgb]{0.56,0.35,0.01}{\textit{#1}}}
\newcommand{\RegionMarkerTok}[1]{#1}
\newcommand{\SpecialCharTok}[1]{\textcolor[rgb]{0.81,0.36,0.00}{\textbf{#1}}}
\newcommand{\SpecialStringTok}[1]{\textcolor[rgb]{0.31,0.60,0.02}{#1}}
\newcommand{\StringTok}[1]{\textcolor[rgb]{0.31,0.60,0.02}{#1}}
\newcommand{\VariableTok}[1]{\textcolor[rgb]{0.00,0.00,0.00}{#1}}
\newcommand{\VerbatimStringTok}[1]{\textcolor[rgb]{0.31,0.60,0.02}{#1}}
\newcommand{\WarningTok}[1]{\textcolor[rgb]{0.56,0.35,0.01}{\textbf{\textit{#1}}}}
\usepackage{longtable,booktabs,array}
\usepackage{calc} % for calculating minipage widths
% Correct order of tables after \paragraph or \subparagraph
\usepackage{etoolbox}
\makeatletter
\patchcmd\longtable{\par}{\if@noskipsec\mbox{}\fi\par}{}{}
\makeatother
% Allow footnotes in longtable head/foot
\IfFileExists{footnotehyper.sty}{\usepackage{footnotehyper}}{\usepackage{footnote}}
\makesavenoteenv{longtable}
\setlength{\emergencystretch}{3em} % prevent overfull lines
\providecommand{\tightlist}{%
  \setlength{\itemsep}{0pt}\setlength{\parskip}{0pt}}
\setcounter{secnumdepth}{-\maxdimen} % remove section numbering
\ifLuaTeX
\usepackage[bidi=basic]{babel}
\else
\usepackage[bidi=default]{babel}
\fi
\babelprovide[main,import]{brazilian}
% get rid of language-specific shorthands (see #6817):
\let\LanguageShortHands\languageshorthands
\def\languageshorthands#1{}
% Preâmbulo do Manual do autor (Bruma) — usado pelo pandoc via --include-in-header
\usepackage{fontspec}
\setmainfont{DejaVu Sans}[Scale=0.92]
\setmonofont{DejaVu Sans Mono}[Scale=0.84]
\newfontfamily\titulofonte{SpaceGrotesk-SemiBold.ttf}[Path=./fonts/]
\usepackage{xcolor}
\definecolor{ciano}{HTML}{0A8F78}
\definecolor{cianoclaro}{HTML}{4FE3C1}
\definecolor{ambar}{HTML}{B26B00}
\definecolor{violeta}{HTML}{6E5BD8}
\definecolor{tinta}{HTML}{0C1614}
\definecolor{fundo}{HTML}{F3F7F6}
\definecolor{shadecolor}{HTML}{F3F7F6}
\usepackage{graphicx}
\graphicspath{{img/}}
\usepackage[a4paper,top=2.3cm,bottom=2.3cm,left=2.3cm,right=2.3cm,headheight=22pt]{geometry}
\usepackage{xurl}
\usepackage{microtype}
\usepackage{titlesec}
\titleformat{\section}{\titulofonte\Large\color{tinta}}{}{0pt}{}[{\color{cianoclaro}\titlerule[1.2pt]}]
\titleformat{\subsection}{\titulofonte\large\color{ciano}}{}{0pt}{}
\titleformat{\subsubsection}{\titulofonte\normalsize\color{tinta}}{}{0pt}{}
\titlespacing*{\section}{0pt}{20pt}{10pt}
\titlespacing*{\subsection}{0pt}{14pt}{6pt}
\usepackage{fancyhdr}
\pagestyle{fancy}
\fancyhf{}
\fancyhead[L]{\raisebox{-0.25\height}{\includegraphics[height=14pt]{bruma-256.png}}\hspace{6pt}{\titulofonte\small Manual do autor}}
\fancyhead[R]{\small\color{gray}Bruma e brunohdmacedo.tech}
\fancyfoot[L]{\small\color{gray}Bruno Henrique Dourado Macedo}
\fancyfoot[R]{\small\thepage}
\renewcommand{\headrulewidth}{0.4pt}
\fancypagestyle{plain}{\fancyhf{}\fancyfoot[R]{\small\thepage}\renewcommand{\headrulewidth}{0pt}}
% blocos de código: quebra de linha e moldura leve
\usepackage{fvextra}
\RecustomVerbatimEnvironment{verbatim}{Verbatim}{breaklines,breakanywhere,fontsize=\small,frame=single,framesep=6pt,rulecolor=\color{black!15}}
\makeatletter\@ifundefined{Highlighting}{}{\RecustomVerbatimEnvironment{Highlighting}{Verbatim}{breaklines,breakanywhere,commandchars=\\\{\},fontsize=\small}}\makeatother
% citações (blockquote) com barra lateral
\usepackage{framed}
\renewenvironment{quote}{\def\FrameCommand{{\color{ambar}\vrule width 3pt}\hspace{8pt}\colorbox{fundo}}\MakeFramed{\advance\hsize-\width\FrameRestore}\small}{\endMakeFramed}
% tabelas
\usepackage{booktabs,longtable,array}
\renewcommand{\arraystretch}{1.25}
\setlength{\LTpre}{6pt}\setlength{\LTpost}{8pt}
% listas mais compactas
\usepackage{enumitem}
\setlist{itemsep=2pt,topsep=4pt}
% capa
\renewcommand{\maketitle}{%
\begin{titlepage}
\centering
\vspace*{2.6cm}
\includegraphics[width=0.55\linewidth]{bruma-logo-claro.png}\par
\vspace{1.6cm}
{\titulofonte\fontsize{30}{36}\selectfont\color{tinta} Manual do autor\par}
\vspace{0.4cm}
{\Large\color{ciano} Bruma e site brunohdmacedo.tech\par}
\vspace{0.3cm}
{\large Como publicar, usar, manter e atualizar tudo\par}
\vspace{2.2cm}
{\large Bruno Henrique Dourado Macedo\par}
\vspace{0.2cm}
{\small Engenheiro Físico (UNILA) · Mestrando em Física Aplicada (PPGFISA/UNILA) · bolsista CNPq\par}
\vspace{0.2cm}
{\small\url{https://brunohdmacedo.tech} \quad·\quad ORCID 0000-0002-8152-0950\par}
\vfill
{\small\color{gray} Versão 1.0.0 · novembro de 2026\par}
\end{titlepage}}
\ifLuaTeX
  \usepackage{selnolig}  % disable illegal ligatures
\fi
\IfFileExists{bookmark.sty}{\usepackage{bookmark}}{\usepackage{hyperref}}
\IfFileExists{xurl.sty}{\usepackage{xurl}}{} % add URL line breaks if available
\urlstyle{same}
\hypersetup{
  pdftitle={Manual do autor: Bruma e site brunohdmacedo.tech},
  pdfauthor={Bruno Henrique Dourado Macedo},
  pdflang={pt-BR},
  colorlinks=true,
  linkcolor={ciano},
  filecolor={Maroon},
  citecolor={Blue},
  urlcolor={ciano},
  pdfcreator={LaTeX via pandoc}}

\title{Manual do autor: Bruma e site brunohdmacedo.tech}
\author{Bruno Henrique Dourado Macedo}
\date{Versão 1.0.0 · novembro de 2026}

\begin{document}
\maketitle

\renewcommand*\contentsname{Sumário}
{
\hypersetup{linkcolor=tinta}
\setcounter{tocdepth}{2}
\tableofcontents
}
Este manual explica \textbf{como você faz tudo}:

\begin{itemize}
\tightlist
\item
  publicar e manter o site pessoal;
\item
  publicar o Bruma no GitHub, no PyPI e no Windows;
\item
  usar cada versão do programa (web, cluster, janela e Windows);
\item
  manter as citações automáticas;
\item
  lançar novas versões;
\item
  compilar os documentos em LaTeX;
\item
  resolver os problemas mais comuns.
\end{itemize}

\begin{quote}
\textbf{Manual ou tutorial?} O \emph{tutorial}
(\texttt{tutorial.\allowbreak{}html} e
\texttt{bruma-\allowbreak{}tutorial.\allowbreak{}pdf}) ensina a
\textbf{usar o Bruma} para analisar cálculos. Este \emph{manual} ensina
a \textbf{publicar e manter} tudo o que foi criado.
\end{quote}

\hypertarget{o-que-foi-criado}{%
\section{0. O que foi criado}\label{o-que-foi-criado}}

\begin{longtable}[]{@{}
  >{\raggedright\arraybackslash}p{(\columnwidth - 4\tabcolsep) * \real{0.3333}}
  >{\raggedright\arraybackslash}p{(\columnwidth - 4\tabcolsep) * \real{0.3333}}
  >{\raggedright\arraybackslash}p{(\columnwidth - 4\tabcolsep) * \real{0.3333}}@{}}
\toprule\noalign{}
\begin{minipage}[b]{\linewidth}\raggedright
Entrega
\end{minipage} & \begin{minipage}[b]{\linewidth}\raggedright
Onde fica
\end{minipage} & \begin{minipage}[b]{\linewidth}\raggedright
Para quê
\end{minipage} \\
\midrule\noalign{}
\endhead
\bottomrule\noalign{}
\endlastfoot
\textbf{Site pessoal} & repositório
\texttt{brunohdmacedo/\allowbreak{}brunohdmacedo.\allowbreak{}github.\allowbreak{}io}
→ brunohdmacedo.tech & portfólio, publicações com citações automáticas,
Bruma em \texttt{/\allowbreak{}bruma/\allowbreak{}} \\
\textbf{Bruma (web)} &
\texttt{bruma-\allowbreak{}dft/\allowbreak{}index.\allowbreak{}html},
publicado em brunohdmacedo.tech/bruma & análise no navegador, sem
instalar nada \\
\textbf{Bruma (linha de comando)} &
\texttt{bruma-\allowbreak{}dft/\allowbreak{}bruma.\allowbreak{}py} &
cluster, scripts, filas SLURM/PBS \\
\textbf{Bruma (janela)} &
\texttt{bruma-\allowbreak{}dft/\allowbreak{}bruma\_\allowbreak{}gui.\allowbreak{}py}
& uso local com cliques (Linux, macOS, Windows) \\
\textbf{Bruma (Windows)} & GitHub → Releases:
\texttt{Bruma-\allowbreak{}Setup-\allowbreak{}<versão>.\allowbreak{}exe}
& quem não usa Python \\
\textbf{Tutorial} & \texttt{tutorial.\allowbreak{}html} e
\texttt{bruma-\allowbreak{}tutorial.\allowbreak{}pdf} (fonte
\texttt{.\allowbreak{}tex}) & como analisar cálculos com o Bruma \\
\textbf{Citações} & \texttt{bruma.\allowbreak{}bib},
\texttt{CITATION.\allowbreak{}cff} & citar o Bruma, o TCC, o artigo e a
base de dados \\
\textbf{Logo} &
\texttt{bruma-\allowbreak{}dft/\allowbreak{}assets/\allowbreak{}} (e
\texttt{logo-\allowbreak{}bruma.\allowbreak{}zip}) & SVG, PNG de 16 a
1024 px, \texttt{.\allowbreak{}ico} \\
\textbf{Citações automáticas do site} &
\texttt{.\allowbreak{}github/\allowbreak{}workflows/\allowbreak{}citacoes.\allowbreak{}yml},
\texttt{scripts/\allowbreak{}atualizar\_\allowbreak{}citacoes.\allowbreak{}py}
& números do OpenAlex e do Google Acadêmico toda semana \\
\textbf{Compilação para Windows} &
\texttt{.\allowbreak{}github/\allowbreak{}workflows/\allowbreak{}windows.\allowbreak{}yml},
\texttt{packaging/\allowbreak{}windows/\allowbreak{}} & gera os
\texttt{.\allowbreak{}exe} e o instalador \\
\textbf{Publicação no PyPI} &
\texttt{.\allowbreak{}github/\allowbreak{}workflows/\allowbreak{}pypi.\allowbreak{}yml},
\texttt{pyproject.\allowbreak{}toml} &
\texttt{pip install bruma-\allowbreak{}dft} \\
\textbf{Ferramentas} &
\texttt{tools/\allowbreak{}versao.\allowbreak{}py},
\texttt{tools/\allowbreak{}copiar\_\allowbreak{}para\_\allowbreak{}site.\allowbreak{}sh}
& trocar a versão em todos os arquivos; copiar o Bruma para o site \\
\textbf{Exemplos de cluster} &
\texttt{examples/\allowbreak{}slurm\_\allowbreak{}bruma.\allowbreak{}sh},
\texttt{pbs\_\allowbreak{}bruma.\allowbreak{}sh},
\texttt{castep\_\allowbreak{}depois\_\allowbreak{}bruma.\allowbreak{}sh}
& jobs prontos \\
\end{longtable}

São \textbf{dois repositórios} no GitHub:

\begin{enumerate}
\def\labelenumi{\arabic{enumi}.}
\tightlist
\item
  \textbf{\texttt{brunohdmacedo.\allowbreak{}github.\allowbreak{}io}:} o
  site (já existe).
\item
  \textbf{\texttt{bruma-\allowbreak{}dft}:} o software (você vai criar).
  O site recebe uma \emph{cópia} da versão web do Bruma na pasta
  \texttt{/\allowbreak{}bruma/\allowbreak{}}.
\end{enumerate}

\begin{center}\rule{0.5\linewidth}{0.5pt}\end{center}

\hypertarget{antes-de-comeuxe7ar}{%
\section{1. Antes de começar}\label{antes-de-comeuxe7ar}}

\hypertarget{contas}{%
\subsection{Contas}\label{contas}}

\begin{longtable}[]{@{}
  >{\raggedright\arraybackslash}p{(\columnwidth - 4\tabcolsep) * \real{0.3333}}
  >{\raggedright\arraybackslash}p{(\columnwidth - 4\tabcolsep) * \real{0.3333}}
  >{\raggedright\arraybackslash}p{(\columnwidth - 4\tabcolsep) * \real{0.3333}}@{}}
\toprule\noalign{}
\begin{minipage}[b]{\linewidth}\raggedright
Conta
\end{minipage} & \begin{minipage}[b]{\linewidth}\raggedright
Para quê
\end{minipage} & \begin{minipage}[b]{\linewidth}\raggedright
Obrigatória?
\end{minipage} \\
\midrule\noalign{}
\endhead
\bottomrule\noalign{}
\endlastfoot
GitHub (\texttt{brunohdmacedo}) & repositórios, site, Actions, Releases
& sim \\
PyPI (pypi.org) & \texttt{pip install bruma-\allowbreak{}dft} & não \\
Zenodo (zenodo.org) & DOI para cada versão do Bruma & não, mas
recomendado \\
SerpApi (serpapi.com) & números do Google Acadêmico no site & não \\
Overleaf (overleaf.com) & compilar o LaTeX sem instalar nada & não \\
\end{longtable}

\hypertarget{programas-no-seu-computador}{%
\subsection{Programas no seu
computador}\label{programas-no-seu-computador}}

\begin{itemize}
\item
  \textbf{Git:} \url{https://git-scm.com}. Configure uma vez:

\begin{Shaded}
\begin{Highlighting}[]
\FunctionTok{git}\NormalTok{ config }\AttributeTok{{-}{-}global}\NormalTok{ user.name }\StringTok{"Bruno Henrique Dourado Macedo"}
\FunctionTok{git}\NormalTok{ config }\AttributeTok{{-}{-}global}\NormalTok{ user.email }\StringTok{"brunohdmacedo@gmail.com"}
\end{Highlighting}
\end{Shaded}
\item
  \textbf{Python 3.8 ou mais novo:} \url{https://python.org}. No
  Windows, marque \emph{Add python.exe to PATH} na instalação.
\item
  \textbf{Opcional:}

  \begin{itemize}
  \tightlist
  \item
    TeX Live ou MiKTeX, para compilar o LaTeX localmente;
  \item
    Inno Setup 6, para gerar o instalador no seu Windows;
  \item
    GitHub Desktop, se preferir não usar a linha de comando.
  \end{itemize}
\end{itemize}

\hypertarget{arquivos-entregues}{%
\subsection{Arquivos entregues}\label{arquivos-entregues}}

\begin{longtable}[]{@{}
  >{\raggedright\arraybackslash}p{(\columnwidth - 2\tabcolsep) * \real{0.5000}}
  >{\raggedright\arraybackslash}p{(\columnwidth - 2\tabcolsep) * \real{0.5000}}@{}}
\toprule\noalign{}
\begin{minipage}[b]{\linewidth}\raggedright
Arquivo
\end{minipage} & \begin{minipage}[b]{\linewidth}\raggedright
Conteúdo
\end{minipage} \\
\midrule\noalign{}
\endhead
\bottomrule\noalign{}
\endlastfoot
\texttt{bruma-\allowbreak{}dft.\allowbreak{}zip} & repositório completo
do Bruma \\
\texttt{brunohdmacedo.\allowbreak{}github.\allowbreak{}io.\allowbreak{}zip}
& site completo, com as páginas antigas mantidas \\
\texttt{logo-\allowbreak{}bruma.\allowbreak{}zip} & todos os arquivos da
logo \\
\end{longtable}

\begin{center}\rule{0.5\linewidth}{0.5pt}\end{center}

\hypertarget{parte-a-site-pessoal-brunohdmacedo.tech}{%
\section{Parte A: site pessoal
(brunohdmacedo.tech)}\label{parte-a-site-pessoal-brunohdmacedo.tech}}

\hypertarget{a1.-publicar-a-versuxe3o-nova-do-site}{%
\subsection{A1. Publicar a versão nova do
site}\label{a1.-publicar-a-versuxe3o-nova-do-site}}

\begin{enumerate}
\def\labelenumi{\arabic{enumi}.}
\item
  Baixe o seu repositório (se ainda não tiver):

\begin{Shaded}
\begin{Highlighting}[]
\FunctionTok{git}\NormalTok{ clone https://github.com/brunohdmacedo/brunohdmacedo.github.io.git}
\BuiltInTok{cd}\NormalTok{ brunohdmacedo.github.io}
\end{Highlighting}
\end{Shaded}
\item
  Descompacte o
  \texttt{brunohdmacedo.\allowbreak{}github.\allowbreak{}io.\allowbreak{}zip}
  \textbf{por cima} dessa pasta. O zip não contém a pasta
  \texttt{.\allowbreak{}git}, então o histórico é preservado.
\item
  Confira o que mudou:

\begin{Shaded}
\begin{Highlighting}[]
\FunctionTok{git}\NormalTok{ status}
\FunctionTok{git}\NormalTok{ diff }\AttributeTok{{-}{-}stat}
\end{Highlighting}
\end{Shaded}
\item
  Envie:

\begin{Shaded}
\begin{Highlighting}[]
\FunctionTok{git}\NormalTok{ add }\AttributeTok{{-}A}
\FunctionTok{git}\NormalTok{ commit }\AttributeTok{{-}m} \StringTok{"Site: Bruma, publicações, citações automáticas"}
\FunctionTok{git}\NormalTok{ push}
\end{Highlighting}
\end{Shaded}
\item
  Em 1 a 3 minutos o site está no ar. Para acompanhar, abra a aba
  \textbf{Actions} do repositório.
\end{enumerate}

\hypertarget{a2.-configurauxe7uxe3o-do-github-pages-uma-vez}{%
\subsection{A2. Configuração do GitHub Pages (uma
vez)}\label{a2.-configurauxe7uxe3o-do-github-pages-uma-vez}}

Em \textbf{Settings → Pages}:

\begin{itemize}
\tightlist
\item
  \textbf{Source:} o repositório tem um workflow
  \texttt{static.\allowbreak{}yml}, que publica o site a cada
  \texttt{push}. Deixe \textbf{Source = GitHub Actions}. Se preferir
  \textbf{Deploy from a branch} (\texttt{main}, pasta
  \texttt{/\allowbreak{}root}), também funciona; nesse caso o
  \texttt{static.\allowbreak{}yml} pode ser apagado.
\item
  \textbf{Custom domain:} \texttt{brunohdmacedo.\allowbreak{}tech}. O
  arquivo \texttt{CNAME} já existe e não deve ser apagado.
\item
  \textbf{Enforce HTTPS:} marcado.
\end{itemize}

O repositório também tem vários workflows de exemplo (Jekyll, Gatsby,
Next.js, npm). Se aparecerem falhando na aba Actions, desative-os em
\textbf{Actions → (workflow) → ··· → Disable workflow}, ou apague os
arquivos de
\texttt{.\allowbreak{}github/\allowbreak{}workflows/\allowbreak{}} que
não usa. \textbf{Não apague} \texttt{static.\allowbreak{}yml} (se usar
Source = GitHub Actions) nem \texttt{citacoes.\allowbreak{}yml}.

\hypertarget{a3.-editar-o-conteuxfado-do-site}{%
\subsection{A3. Editar o conteúdo do
site}\label{a3.-editar-o-conteuxfado-do-site}}

Tudo está no \texttt{index.\allowbreak{}html} da raiz.

\begin{longtable}[]{@{}
  >{\raggedright\arraybackslash}p{(\columnwidth - 4\tabcolsep) * \real{0.3333}}
  >{\raggedright\arraybackslash}p{(\columnwidth - 4\tabcolsep) * \real{0.3333}}
  >{\raggedright\arraybackslash}p{(\columnwidth - 4\tabcolsep) * \real{0.3333}}@{}}
\toprule\noalign{}
\begin{minipage}[b]{\linewidth}\raggedright
Quero mudar\ldots{}
\end{minipage} & \begin{minipage}[b]{\linewidth}\raggedright
Onde
\end{minipage} & \begin{minipage}[b]{\linewidth}\raggedright
Como
\end{minipage} \\
\midrule\noalign{}
\endhead
\bottomrule\noalign{}
\endlastfoot
Contadores do topo (4 anos, 3 anos, 7 prêmios, 20 apresentações) &
\texttt{<div class="stats">} & troque o \texttt{data-\allowbreak{}n="…"}
(número) e o texto do \texttt{<span>} \\
Texto digitado no topo &
\texttt{site/\allowbreak{}js/\allowbreak{}main.\allowbreak{}js}, lista
\texttt{var ph=[…]} & edite as frases entre aspas \\
Texto ``sobre'' & \texttt{<section id="sobre">} & edite os parágrafos
\texttt{<p>} \\
Linha do tempo & \texttt{<section id="trajetoria">}, blocos
\texttt{<div class="ti">} & copie um bloco e edite;
\texttt{class="ti m"} destaca o ponto em âmbar \\
Prêmios & \texttt{<section id="premios">} & copie uma linha e edite \\
Publicações & \texttt{<div id="pubs">} & veja a seção A4 \\
Cartão do Bruma & comentário
\texttt{<!-\allowbreak{}-\allowbreak{} Bruma:\allowbreak{}} na seção
laboratório & texto e botões \\
\end{longtable}

Para testar localmente antes de enviar:

\begin{Shaded}
\begin{Highlighting}[]
\ExtensionTok{python3} \AttributeTok{{-}m}\NormalTok{ http.server 8000}
\end{Highlighting}
\end{Shaded}

Depois abra
\texttt{http:\allowbreak{}/\allowbreak{}/\allowbreak{}localhost:\allowbreak{}8000}
no navegador.

\hypertarget{a4.-adicionar-uma-publicauxe7uxe3o}{%
\subsection{A4. Adicionar uma
publicação}\label{a4.-adicionar-uma-publicauxe7uxe3o}}

Cada publicação é um bloco dentro de \texttt{<div id="pubs">}:

\begin{Shaded}
\begin{Highlighting}[]
\DataTypeTok{\textless{}}\KeywordTok{div} \ErrorTok{class}\OtherTok{=}\StringTok{"pub"} \ErrorTok{data{-}t}\OtherTok{=}\StringTok{"art"} \ErrorTok{data{-}title}\OtherTok{=}\StringTok{"Título exato da publicação"} \ErrorTok{data{-}doi}\OtherTok{=}\StringTok{"10.xxxx/yyyy"}\DataTypeTok{\textgreater{}}
  \DataTypeTok{\textless{}}\KeywordTok{time}\DataTypeTok{\textgreater{}}\NormalTok{2026}\DataTypeTok{\textless{}/}\KeywordTok{time}\DataTypeTok{\textgreater{}}
  \DataTypeTok{\textless{}}\KeywordTok{div}\DataTypeTok{\textgreater{}}
    \DataTypeTok{\textless{}}\KeywordTok{b}\DataTypeTok{\textgreater{}\textless{}}\KeywordTok{a} \ErrorTok{href}\OtherTok{=}\StringTok{"https://link.da.publicacao"} \ErrorTok{target}\OtherTok{=}\StringTok{"\_blank"} \ErrorTok{rel}\OtherTok{=}\StringTok{"noopener"} \ErrorTok{style}\OtherTok{=}\StringTok{"color:inherit"}\DataTypeTok{\textgreater{}}\NormalTok{Título exato da publicação}\DataTypeTok{\textless{}/}\KeywordTok{a}\DataTypeTok{\textgreater{}\textless{}/}\KeywordTok{b}\DataTypeTok{\textgreater{}}
    \DataTypeTok{\textless{}}\KeywordTok{span}\DataTypeTok{\textgreater{}}\NormalTok{Macedo, B. H. D.; Coautor, X. · Revista, v. 1, p. 1–10, 2026}\DataTypeTok{\textless{}/}\KeywordTok{span}\DataTypeTok{\textgreater{}}
  \DataTypeTok{\textless{}/}\KeywordTok{div}\DataTypeTok{\textgreater{}}
\DataTypeTok{\textless{}/}\KeywordTok{div}\DataTypeTok{\textgreater{}}
\end{Highlighting}
\end{Shaded}

\begin{itemize}
\tightlist
\item
  \textbf{\texttt{data-\allowbreak{}t}} define o filtro: \texttt{art}
  (artigo), \texttt{cong} (congresso), \texttt{tcc}, \texttt{dat} (base
  de dados) ou \texttt{sw} (software). Para criar um filtro novo, copie
  um botão \texttt{<button data-\allowbreak{}f="…">} dentro de
  \texttt{<div class="fl" id="fl">}.
\item
  \textbf{\texttt{data-\allowbreak{}title}} (o título exato) e
  \textbf{\texttt{data-\allowbreak{}doi}} (quando existir) permitem que
  o selo ``citado N×'' apareça sozinho.
\item
  A ordem na lista é a ordem do arquivo. Mantenha as mais recentes no
  topo.
\end{itemize}

\hypertarget{a5.-citauxe7uxf5es-automuxe1ticas-google-acaduxeamico-e-openalex}{%
\subsection{A5. Citações automáticas (Google Acadêmico e
OpenAlex)}\label{a5.-citauxe7uxf5es-automuxe1ticas-google-acaduxeamico-e-openalex}}

\textbf{Como funciona:}

\begin{itemize}
\tightlist
\item
  \textbf{\texttt{citacoes.\allowbreak{}yml}:} roda toda segunda-feira
  às 9h (UTC) e chama o
  \texttt{scripts/\allowbreak{}atualizar\_\allowbreak{}citacoes.\allowbreak{}py}.
\item
  \textbf{\texttt{atualizar\_\allowbreak{}citacoes.\allowbreak{}py}:}
  consulta o \textbf{OpenAlex} pelo seu ORCID (gratuito) e o
  \textbf{Google Acadêmico} pelo SerpApi (se houver chave), e grava o
  resultado em \texttt{data/\allowbreak{}citacoes.\allowbreak{}json}.
\item
  \textbf{\texttt{site/\allowbreak{}js/\allowbreak{}citacoes.\allowbreak{}js}:}
  lê o JSON direto do GitHub e mostra o total, o índice h e os selos em
  cada publicação. Se o JSON estiver vazio, consulta o OpenAlex ao vivo.
\end{itemize}

\textbf{Ativar (uma vez):}

\begin{enumerate}
\def\labelenumi{\arabic{enumi}.}
\tightlist
\item
  Faça o push do site (A1).
\item
  \textbf{Actions → Atualizar citações → Run workflow}.
\item
  Para os números do Google Acadêmico:

  \begin{itemize}
  \tightlist
  \item
    crie uma conta em \url{https://serpapi.com} (o plano gratuito basta
    para uma busca por semana);
  \item
    copie a API key;
  \item
    em \textbf{Settings → Secrets and variables → Actions → New
    repository secret}, use o nome \texttt{SERPAPI\_\allowbreak{}KEY} e
    cole a chave como valor;
  \item
    rode o workflow de novo.
  \end{itemize}
\end{enumerate}

Se o workflow falhar com \emph{permission denied} ao fazer push, vá em
\textbf{Settings → Actions → General → Workflow permissions} e marque
\textbf{Read and write permissions}.

\begin{quote}
O Google Acadêmico proíbe raspagem automática, por isso o caminho é o
SerpApi. O OpenAlex costuma mostrar menos citações que o Scholar, porque
indexa menos teses e anais.
\end{quote}

\hypertarget{a6.-google-e-buscadores}{%
\subsection{A6. Google e buscadores}\label{a6.-google-e-buscadores}}

\begin{itemize}
\tightlist
\item
  O \texttt{sitemap.\allowbreak{}xml} e o
  \texttt{robots.\allowbreak{}txt} já estão na raiz.
\item
  No \textbf{Google Search Console}
  (\url{https://search.google.com/search-console}), o arquivo de
  verificação \texttt{googleea98cecbf58918b7.\allowbreak{}html} já está
  no site. Envie
  \texttt{https:\allowbreak{}/\allowbreak{}/\allowbreak{}brunohdmacedo.\allowbreak{}tech/\allowbreak{}sitemap.\allowbreak{}xml}
  em \textbf{Sitemaps}.
\item
  O endereço antigo \texttt{/\allowbreak{}castep/\allowbreak{}}
  redireciona para \texttt{/\allowbreak{}bruma/\allowbreak{}}. Não
  apague a pasta \texttt{castep/\allowbreak{}}.
\end{itemize}

\begin{center}\rule{0.5\linewidth}{0.5pt}\end{center}

\hypertarget{parte-b-reposituxf3rio-do-bruma-bruma-dft}{%
\section{\texorpdfstring{Parte B: repositório do Bruma
(\texttt{bruma-\allowbreak{}dft})}{Parte B: repositório do Bruma ()}}\label{parte-b-reposituxf3rio-do-bruma-bruma-dft}}

\hypertarget{b1.-criar-e-enviar-uma-vez}{%
\subsection{B1. Criar e enviar (uma
vez)}\label{b1.-criar-e-enviar-uma-vez}}

\begin{enumerate}
\def\labelenumi{\arabic{enumi}.}
\item
  No GitHub, abra \textbf{New repository}:

  \begin{itemize}
  \tightlist
  \item
    nome: \texttt{bruma-\allowbreak{}dft};
  \item
    visibilidade: \textbf{Public};
  \item
    sem README (o zip já tem).
  \end{itemize}
\item
  No seu computador:

\begin{Shaded}
\begin{Highlighting}[]
\FunctionTok{unzip}\NormalTok{ bruma{-}dft.zip}
\BuiltInTok{cd}\NormalTok{ bruma{-}dft}
\FunctionTok{git}\NormalTok{ init }\AttributeTok{{-}b}\NormalTok{ main}
\FunctionTok{git}\NormalTok{ add }\AttributeTok{{-}A}
\FunctionTok{git}\NormalTok{ commit }\AttributeTok{{-}m} \StringTok{"Bruma 1.0.0"}
\FunctionTok{git}\NormalTok{ remote add origin https://github.com/brunohdmacedo/bruma{-}dft.git}
\FunctionTok{git}\NormalTok{ push }\AttributeTok{{-}u}\NormalTok{ origin main}
\end{Highlighting}
\end{Shaded}
\item
  \textbf{Licença:} \textbf{Add file → Create new file}, com o nome
  \texttt{LICENSE}. Clique em \textbf{Choose a license template} e
  escolha \textbf{MIT} (comum em software científico).
\item
  \textbf{Descrição e tópicos} (engrenagem ao lado de \emph{About}):

  \begin{itemize}
  \tightlist
  \item
    descrição: \emph{Bandas, DOS, PDOS, densidade de carga, defeitos e
    termodinâmica a partir de cálculos CASTEP};
  \item
    site:
    \texttt{https:\allowbreak{}/\allowbreak{}/\allowbreak{}brunohdmacedo.\allowbreak{}tech/\allowbreak{}bruma/\allowbreak{}};
  \item
    tópicos: \texttt{castep}, \texttt{dft},
    \texttt{band-\allowbreak{}structure},
    \texttt{density-\allowbreak{}of-\allowbreak{}states},
    \texttt{charge-\allowbreak{}density},
    \texttt{materials-\allowbreak{}science},
    \texttt{hydrogen-\allowbreak{}storage}.
  \end{itemize}
\end{enumerate}

O arquivo \texttt{CITATION.\allowbreak{}cff} faz aparecer o botão
\textbf{Cite this repository} na página do repositório.

\hypertarget{b2.-versuxe3o-web-no-reposituxf3rio-opcional}{%
\subsection{B2. Versão web no repositório
(opcional)}\label{b2.-versuxe3o-web-no-reposituxf3rio-opcional}}

A versão oficial da web fica no seu site
(\texttt{/\allowbreak{}bruma/\allowbreak{}}). Se quiser também um
endereço no próprio repositório: \textbf{Settings → Pages → Deploy from
a branch → main / (root)}. Ele ficará em
\texttt{https:\allowbreak{}/\allowbreak{}/\allowbreak{}brunohdmacedo.\allowbreak{}github.\allowbreak{}io/\allowbreak{}bruma-\allowbreak{}dft/\allowbreak{}}.

\hypertarget{b3.-atualizar-a-cuxf3pia-do-bruma-no-site}{%
\subsection{B3. Atualizar a cópia do Bruma no
site}\label{b3.-atualizar-a-cuxf3pia-do-bruma-no-site}}

Sempre que mudar o \texttt{index.\allowbreak{}html}, o
\texttt{tutorial.\allowbreak{}html} ou os PDFs do Bruma:

\begin{Shaded}
\begin{Highlighting}[]
\BuiltInTok{cd}\NormalTok{ bruma{-}dft}
\FunctionTok{bash}\NormalTok{ tools/copiar\_para\_site.sh ../brunohdmacedo.github.io}
\BuiltInTok{cd}\NormalTok{ ../brunohdmacedo.github.io}
\FunctionTok{git}\NormalTok{ add bruma}
\FunctionTok{git}\NormalTok{ commit }\AttributeTok{{-}m} \StringTok{"Atualiza o Bruma"}
\FunctionTok{git}\NormalTok{ push}
\end{Highlighting}
\end{Shaded}

No Windows, sem Bash, copie à mão para a pasta
\texttt{bruma/\allowbreak{}} do site estes arquivos:

\begin{itemize}
\tightlist
\item
  \texttt{index.\allowbreak{}html}, \texttt{tutorial.\allowbreak{}html},
  \texttt{bruma-\allowbreak{}tutorial.\allowbreak{}pdf},
  \texttt{bruma-\allowbreak{}tutorial.\allowbreak{}tex},
  \texttt{bruma.\allowbreak{}bib} e \texttt{CITATION.\allowbreak{}cff};
\item
  as pastas \texttt{img/\allowbreak{}} e \texttt{assets/\allowbreak{}}.
\end{itemize}

\begin{center}\rule{0.5\linewidth}{0.5pt}\end{center}

\hypertarget{parte-c-usar-o-bruma-nas-quatro-formas}{%
\section{Parte C: usar o Bruma nas quatro
formas}\label{parte-c-usar-o-bruma-nas-quatro-formas}}

Detalhes de cada análise estão no \textbf{tutorial}. Aqui está como pôr
cada forma para funcionar.

\hypertarget{c1.-web}{%
\subsection{C1. Web}\label{c1.-web}}

Abra \url{https://brunohdmacedo.tech/bruma/} e arraste os arquivos do
cálculo (\texttt{.\allowbreak{}bands},
\texttt{.\allowbreak{}pdos\_\allowbreak{}bin},
\texttt{.\allowbreak{}castep}, \texttt{.\allowbreak{}cell},
\texttt{.\allowbreak{}den\_\allowbreak{}fmt} ou \texttt{CHGCAR}). Os
arquivos \textbf{não saem do seu computador}: tudo roda no navegador. A
aba \textbf{Sobre} e a seção \textbf{Sobre o Bruma} trazem o autor e a
citação.

\hypertarget{c2.-linha-de-comando-no-cluster}{%
\subsection{C2. Linha de comando no
cluster}\label{c2.-linha-de-comando-no-cluster}}

\textbf{Opção 1, sem instalar:} copie o script e use o Python do
cluster.

\begin{Shaded}
\begin{Highlighting}[]
\FunctionTok{scp}\NormalTok{ bruma.py usuario@cluster:\textasciitilde{}/                  }\CommentTok{\# do seu computador}
\FunctionTok{ssh}\NormalTok{ usuario@cluster}
\ExtensionTok{module}\NormalTok{ load python                               }\CommentTok{\# se o cluster usar módulos}
\ExtensionTok{python3} \AttributeTok{{-}m}\NormalTok{ pip install }\AttributeTok{{-}{-}user}\NormalTok{ numpy matplotlib   }\CommentTok{\# uma vez, se faltar}
\BuiltInTok{cd}\NormalTok{ /caminho/do/calculo}
\ExtensionTok{python3}\NormalTok{ \textasciitilde{}/bruma.py mgh2\_Fe }\AttributeTok{{-}e}\NormalTok{ Mg,H,Fe}
\end{Highlighting}
\end{Shaded}

\textbf{Opção 2, instalado:} quando o Bruma estiver no PyPI (Parte E).

\begin{Shaded}
\begin{Highlighting}[]
\ExtensionTok{python3} \AttributeTok{{-}m}\NormalTok{ pip install }\AttributeTok{{-}{-}user}\NormalTok{ bruma{-}dft}
\BuiltInTok{export} \VariableTok{PATH}\OperatorTok{=}\StringTok{"}\VariableTok{$HOME}\StringTok{/.local/bin:}\VariableTok{$PATH}\StringTok{"}             \CommentTok{\# ponha no \textasciitilde{}/.bashrc}
\ExtensionTok{bruma}\NormalTok{ mgh2\_Fe }\AttributeTok{{-}e}\NormalTok{ Mg,H,Fe}
\end{Highlighting}
\end{Shaded}

\textbf{Em fila (SLURM/PBS):} use os exemplos da pasta
\texttt{examples/\allowbreak{}}.

\begin{Shaded}
\begin{Highlighting}[]
\ExtensionTok{sbatch}\NormalTok{ slurm\_bruma.sh mgh2\_Fe                    }\CommentTok{\# SLURM}
\ExtensionTok{qsub} \AttributeTok{{-}v}\NormalTok{ SEED=mgh2\_Fe pbs\_bruma.sh                }\CommentTok{\# PBS/Torque}
\end{Highlighting}
\end{Shaded}

O
\texttt{castep\_\allowbreak{}depois\_\allowbreak{}bruma.\allowbreak{}sh}
roda o CASTEP e o Bruma em sequência no mesmo job.

\textbf{Trazer os resultados} (no seu computador):

\begin{Shaded}
\begin{Highlighting}[]
\FunctionTok{rsync} \AttributeTok{{-}av}\NormalTok{ usuario@cluster:/caminho/do/calculo/mgh2\_Fe\_plots ./}
\end{Highlighting}
\end{Shaded}

\textbf{Comandos úteis:}

\begin{longtable}[]{@{}
  >{\raggedright\arraybackslash}p{(\columnwidth - 2\tabcolsep) * \real{0.5000}}
  >{\raggedright\arraybackslash}p{(\columnwidth - 2\tabcolsep) * \real{0.5000}}@{}}
\toprule\noalign{}
\begin{minipage}[b]{\linewidth}\raggedright
Comando
\end{minipage} & \begin{minipage}[b]{\linewidth}\raggedright
Faz
\end{minipage} \\
\midrule\noalign{}
\endhead
\bottomrule\noalign{}
\endlastfoot
\texttt{bruma mgh2\_\allowbreak{}Fe} & análise completa (bandas, DOS,
PDOS, análise, defeito, densidade se houver) \\
\texttt{bruma mgh2\_\allowbreak{}Fe -\allowbreak{}-\allowbreak{}ref mgh2\_\allowbreak{}puro}
& defeito com a estrutura pura de referência \\
\texttt{bruma mgh2\_\allowbreak{}Fe -\allowbreak{}-\allowbreak{}gerar-\allowbreak{}referencias}
& gera os \texttt{.\allowbreak{}cell}/\texttt{.\allowbreak{}param} das
referências \\
\texttt{bruma mgh2\_\allowbreak{}Fe -\allowbreak{}-\allowbreak{}planes todos -\allowbreak{}-\allowbreak{}vesta -\allowbreak{}-\allowbreak{}subtract mgh2\_\allowbreak{}puro.\allowbreak{}den\_\allowbreak{}fmt -\allowbreak{}-\allowbreak{}diff-\allowbreak{}sign puro-\allowbreak{}defeito}
& densidade, galeria de planos e pacote VESTA com Δρ \\
\texttt{bruma -\allowbreak{}-\allowbreak{}sobre} & autor, links e como
citar \\
\texttt{bruma -\allowbreak{}h} & todas as opções \\
\end{longtable}

Sem o pacote instalado, troque \texttt{bruma} por
\texttt{python3 bruma.\allowbreak{}py}.

\hypertarget{c3.-interface-gruxe1fica}{%
\subsection{C3. Interface gráfica}\label{c3.-interface-gruxe1fica}}

\textbf{Instalar e abrir:}

\begin{Shaded}
\begin{Highlighting}[]
\ExtensionTok{pip}\NormalTok{ install }\StringTok{"bruma{-}dft[gui]"}                     \CommentTok{\# ou: pip install numpy matplotlib pillow}
\ExtensionTok{bruma{-}gui}                                        \CommentTok{\# ou: python3 bruma\_gui.py}
\end{Highlighting}
\end{Shaded}

\begin{itemize}
\tightlist
\item
  \textbf{Linux:} se aparecer \emph{No module named tkinter}, instale o
  \texttt{python3-\allowbreak{}tk}
  (\texttt{sudo apt install python3-\allowbreak{}tk}).
\item
  \textbf{macOS e Windows:} o Python de python.org já inclui o Tkinter.
\end{itemize}

\textbf{Na janela:}

\begin{enumerate}
\def\labelenumi{\arabic{enumi}.}
\tightlist
\item
  \textbf{Análise:} escolha a pasta, o cálculo e as etapas, e clique em
  \textbf{Executar análise}. O botão \textbf{Ver comando} mostra o
  comando equivalente para o cluster.
\item
  \textbf{Resultados:} lista e prévia das figuras. Um clique duplo abre
  o arquivo (Origin, VESTA, editor).
\item
  \textbf{Sobre:} autor, foto, redes e citação com botões de copiar.
\end{enumerate}

\hypertarget{c4.-windows-instalador}{%
\subsection{C4. Windows (instalador)}\label{c4.-windows-instalador}}

\begin{enumerate}
\def\labelenumi{\arabic{enumi}.}
\tightlist
\item
  Baixe o
  \texttt{Bruma-\allowbreak{}Setup-\allowbreak{}<versão>.\allowbreak{}exe}
  em \url{https://github.com/brunohdmacedo/bruma-dft/releases/latest}.
\item
  Execute. Se aparecer \textbf{``O Windows protegeu o computador''} (o
  programa não é assinado digitalmente), clique em \textbf{Mais
  informações → Executar assim mesmo}.
\item
  Escolha os atalhos e, se quiser, a opção de colocar o comando
  \texttt{bruma} no PATH.
\item
  Abra o \textbf{Bruma} pelo menu Iniciar. A linha de comando fica em
  \texttt{C:\allowbreak{}\textbackslash\allowbreak{}Program Files\textbackslash\allowbreak{}Bruma\textbackslash\allowbreak{}cli\textbackslash\allowbreak{}bruma.\allowbreak{}exe}.
\end{enumerate}

\begin{center}\rule{0.5\linewidth}{0.5pt}\end{center}

\hypertarget{parte-d-programa-para-windows}{%
\section{Parte D: programa para
Windows}\label{parte-d-programa-para-windows}}

Um \texttt{.\allowbreak{}exe} de Windows precisa ser compilado
\textbf{num Windows}. Há três caminhos.

\hypertarget{d1.-pelo-github-actions-com-uma-versuxe3o-publicada-recomendado}{%
\subsection{D1. Pelo GitHub Actions, com uma versão publicada
(recomendado)}\label{d1.-pelo-github-actions-com-uma-versuxe3o-publicada-recomendado}}

\begin{Shaded}
\begin{Highlighting}[]
\FunctionTok{git}\NormalTok{ tag v1.0.0}
\FunctionTok{git}\NormalTok{ push }\AttributeTok{{-}{-}tags}
\end{Highlighting}
\end{Shaded}

O workflow \texttt{windows.\allowbreak{}yml} roda numa máquina Windows
do GitHub:

\begin{enumerate}
\def\labelenumi{\arabic{enumi}.}
\tightlist
\item
  instala o Python e as dependências;
\item
  baixa a sua foto do Lattes e a embute no programa;
\item
  compila o \texttt{Bruma.\allowbreak{}exe} (janela) e o
  \texttt{bruma.\allowbreak{}exe} (linha de comando) com o PyInstaller;
\item
  testa o
  \texttt{bruma.\allowbreak{}exe -\allowbreak{}-\allowbreak{}versao} e o
  \texttt{bruma.\allowbreak{}exe -\allowbreak{}-\allowbreak{}sobre};
\item
  monta o
  \texttt{Bruma-\allowbreak{}Setup-\allowbreak{}<versão>.\allowbreak{}exe}
  com o Inno Setup;
\item
  publica os três arquivos em \textbf{Releases}.
\end{enumerate}

Acompanhe em \textbf{Actions → Windows (.exe e instalador)}. Leva cerca
de 5 a 10 minutos.

\hypertarget{d2.-pelo-github-actions-sem-publicar-versuxe3o-para-testar}{%
\subsection{D2. Pelo GitHub Actions, sem publicar versão (para
testar)}\label{d2.-pelo-github-actions-sem-publicar-versuxe3o-para-testar}}

Em \textbf{Actions → Windows (.exe e instalador) → Run workflow}. Ao
terminar, baixe o zip \textbf{Bruma-Windows} em \emph{Artifacts}, na
página da execução.

\hypertarget{d3.-no-seu-windows}{%
\subsection{D3. No seu Windows}\label{d3.-no-seu-windows}}

\begin{enumerate}
\def\labelenumi{\arabic{enumi}.}
\tightlist
\item
  Instale o Python e, para gerar o instalador, o Inno Setup 6
  (\url{https://jrsoftware.org/isdl.php}).
\item
  Dê dois cliques em
  \texttt{packaging\textbackslash\allowbreak{}windows\textbackslash\allowbreak{}build\_\allowbreak{}windows.\allowbreak{}bat}.
\item
  Os arquivos aparecem em \texttt{dist\textbackslash\allowbreak{}}.
\end{enumerate}

\hypertarget{d4.-assinatura-digital-opcional}{%
\subsection{D4. Assinatura digital
(opcional)}\label{d4.-assinatura-digital-opcional}}

O aviso do SmartScreen some quando o programa é assinado com um
certificado de \emph{code signing}. Há serviços gratuitos para projetos
de código aberto, como o SignPath Foundation; certificados comerciais
são pagos. Sem assinatura, o programa funciona normalmente, com o aviso
descrito em C4.

\begin{center}\rule{0.5\linewidth}{0.5pt}\end{center}

\hypertarget{parte-e-pypi-e-doi-no-zenodo}{%
\section{Parte E: PyPI e DOI no
Zenodo}\label{parte-e-pypi-e-doi-no-zenodo}}

\hypertarget{e1.-publicar-no-pypi-para-pip-install-bruma-dft}{%
\subsection{\texorpdfstring{E1. Publicar no PyPI (para
\texttt{pip install bruma-\allowbreak{}dft})}{E1. Publicar no PyPI (para )}}\label{e1.-publicar-no-pypi-para-pip-install-bruma-dft}}

\textbf{Automático, recomendado (Trusted Publishing, sem senha):}

\begin{enumerate}
\def\labelenumi{\arabic{enumi}.}
\tightlist
\item
  Crie a conta em \url{https://pypi.org} e ative a verificação em duas
  etapas.
\item
  Em \textbf{Your account → Publishing → Add a new pending publisher},
  preencha:

  \begin{itemize}
  \tightlist
  \item
    PyPI project name: \texttt{bruma-\allowbreak{}dft};
  \item
    Owner: \texttt{brunohdmacedo};
  \item
    Repository: \texttt{bruma-\allowbreak{}dft};
  \item
    Workflow: \texttt{pypi.\allowbreak{}yml};
  \item
    Environment: \texttt{pypi}.
  \end{itemize}
\item
  No GitHub, crie o ambiente em \textbf{Settings → Environments → New
  environment}, com o nome \texttt{pypi}.
\item
  A cada \textbf{Release} publicada (D1), o workflow
  \texttt{pypi.\allowbreak{}yml} envia a versão ao PyPI.
\end{enumerate}

\textbf{Manual:}

\begin{Shaded}
\begin{Highlighting}[]
\ExtensionTok{python} \AttributeTok{{-}m}\NormalTok{ pip install }\AttributeTok{{-}{-}upgrade}\NormalTok{ build twine}
\ExtensionTok{python} \AttributeTok{{-}m}\NormalTok{ build                  }\CommentTok{\# cria dist/*.whl e dist/*.tar.gz}
\ExtensionTok{python} \AttributeTok{{-}m}\NormalTok{ twine upload dist/}\PreprocessorTok{*}    \CommentTok{\# usuário: \_\_token\_\_ ; senha: o token do PyPI}
\end{Highlighting}
\end{Shaded}

Para ensaiar antes, publique no TestPyPI
(\texttt{twine upload -\allowbreak{}r testpypi dist/\allowbreak{}*}).

\hypertarget{e2.-doi-para-cada-versuxe3o-zenodo}{%
\subsection{E2. DOI para cada versão
(Zenodo)}\label{e2.-doi-para-cada-versuxe3o-zenodo}}

Um DOI torna o Bruma citável como software em artigos, no TCC e na
dissertação.

\begin{enumerate}
\def\labelenumi{\arabic{enumi}.}
\tightlist
\item
  Entre em \url{https://zenodo.org} com a conta do GitHub.
\item
  Em \textbf{(seu nome) → GitHub}, ative o repositório
  \texttt{bruma-\allowbreak{}dft}.
\item
  Publique uma Release (D1). O Zenodo arquiva a versão e gera um DOI.
\item
  Copie o DOI. Acrescente
  \texttt{doi:\allowbreak{} "10.\allowbreak{}5281/\allowbreak{}zenodo.\allowbreak{}XXXXXXX"}
  ao \texttt{CITATION.\allowbreak{}cff},
  \texttt{doi = \{10.\allowbreak{}5281/\allowbreak{}zenodo.\allowbreak{}XXXXXXX\}}
  à entrada \texttt{macedo2026bruma} do \texttt{bruma.\allowbreak{}bib},
  e o selo de DOI do Zenodo ao topo do \texttt{README.\allowbreak{}md}.
\end{enumerate}

O Zenodo dá um DOI por versão e também um \emph{concept DOI}, que aponta
sempre para a mais recente. Para citar o software em geral, use o
concept DOI.

\begin{center}\rule{0.5\linewidth}{0.5pt}\end{center}

\hypertarget{parte-f-documentos-em-latex}{%
\section{Parte F: documentos em
LaTeX}\label{parte-f-documentos-em-latex}}

\begin{longtable}[]{@{}ll@{}}
\toprule\noalign{}
Documento & Arquivos \\
\midrule\noalign{}
\endhead
\bottomrule\noalign{}
\endlastfoot
Tutorial do Bruma &
\texttt{bruma-\allowbreak{}tutorial.\allowbreak{}tex} +
\texttt{bruma.\allowbreak{}bib} + \texttt{img/\allowbreak{}} \\
Este manual &
\texttt{docs/\allowbreak{}manual/\allowbreak{}manual-\allowbreak{}bruma.\allowbreak{}tex}
+ \texttt{docs/\allowbreak{}manual/\allowbreak{}img/\allowbreak{}} \\
\end{longtable}

\hypertarget{f1.-compilar-no-seu-computador}{%
\subsection{F1. Compilar no seu
computador}\label{f1.-compilar-no-seu-computador}}

Os dois documentos usam \textbf{XeLaTeX} (por causa do português com o
\texttt{polyglossia}):

\begin{Shaded}
\begin{Highlighting}[]
\ExtensionTok{xelatex}\NormalTok{ bruma{-}tutorial}
\ExtensionTok{bibtex}\NormalTok{  bruma{-}tutorial}
\ExtensionTok{xelatex}\NormalTok{ bruma{-}tutorial}
\ExtensionTok{xelatex}\NormalTok{ bruma{-}tutorial}
\end{Highlighting}
\end{Shaded}

Para o manual, que não usa BibTeX:

\begin{Shaded}
\begin{Highlighting}[]
\BuiltInTok{cd}\NormalTok{ docs/manual}
\ExtensionTok{xelatex}\NormalTok{ manual{-}bruma}
\ExtensionTok{xelatex}\NormalTok{ manual{-}bruma}
\end{Highlighting}
\end{Shaded}

\hypertarget{f2.-compilar-no-overleaf}{%
\subsection{F2. Compilar no Overleaf}\label{f2.-compilar-no-overleaf}}

\begin{enumerate}
\def\labelenumi{\arabic{enumi}.}
\tightlist
\item
  \textbf{New Project → Upload Project} e envie um zip com o
  \texttt{.\allowbreak{}tex}, o \texttt{.\allowbreak{}bib} e a pasta
  \texttt{img/\allowbreak{}}.
\item
  \textbf{Menu → Compiler → XeLaTeX}.
\item
  Clique em \textbf{Recompile}.
\end{enumerate}

\hypertarget{f3.-regerar-o-manual-a-partir-do-markdown}{%
\subsection{F3. Regerar o manual a partir do
Markdown}\label{f3.-regerar-o-manual-a-partir-do-markdown}}

A fonte deste manual é o
\texttt{docs/\allowbreak{}MANUAL.\allowbreak{}md} (a versão do GitHub).
Para atualizar o LaTeX e o PDF depois de editá-lo:

\begin{Shaded}
\begin{Highlighting}[]
\BuiltInTok{cd}\NormalTok{ docs/manual}
\FunctionTok{bash}\NormalTok{ gerar\_manual.sh             }\CommentTok{\# pandoc + xelatex}
\end{Highlighting}
\end{Shaded}

\hypertarget{f4.-citar-no-tcc-e-na-dissertauxe7uxe3o}{%
\subsection{F4. Citar no TCC e na
dissertação}\label{f4.-citar-no-tcc-e-na-dissertauxe7uxe3o}}

Inclua as entradas do \texttt{bruma.\allowbreak{}bib} no seu
\texttt{.\allowbreak{}bib} e use
\texttt{\textbackslash\allowbreak{}cite\{macedo2026bruma\}} (software),
\texttt{\textbackslash\allowbreak{}cite\{macedo2025tcc\}} (TCC),
\texttt{\textbackslash\allowbreak{}cite\{macedo2024rbic\}} (artigo) e
\texttt{\textbackslash\allowbreak{}cite\{macedo2024dataset\}} (base de
dados). Com o \textbf{abnTeX2}, as mesmas entradas funcionam; para o
TCC, troque \texttt{@misc} por \texttt{@monography} com
\texttt{type = \{Trabalho de Conclusão de Curso\}}, se o seu modelo
pedir.

Antes de usar, confira os campos marcados com \textbf{CONFIRA} nos
comentários do \texttt{bruma.\allowbreak{}bib}: o título e o ano da
dissertação, e o nome completo do coautor do trabalho de 2025.

\begin{center}\rule{0.5\linewidth}{0.5pt}\end{center}

\hypertarget{parte-g-logo}{%
\section{Parte G: logo}\label{parte-g-logo}}

\begin{longtable}[]{@{}
  >{\raggedright\arraybackslash}p{(\columnwidth - 2\tabcolsep) * \real{0.5000}}
  >{\raggedright\arraybackslash}p{(\columnwidth - 2\tabcolsep) * \real{0.5000}}@{}}
\toprule\noalign{}
\begin{minipage}[b]{\linewidth}\raggedright
Arquivo
\end{minipage} & \begin{minipage}[b]{\linewidth}\raggedright
Uso
\end{minipage} \\
\midrule\noalign{}
\endhead
\bottomrule\noalign{}
\endlastfoot
\texttt{bruma-\allowbreak{}mark.\allowbreak{}svg} & símbolo (ícones,
avatar, slides) \\
\texttt{bruma-\allowbreak{}mark-\allowbreak{}pequeno.\allowbreak{}svg} &
símbolo simplificado para 16--32 px (favicon) \\
\texttt{bruma-\allowbreak{}logo-\allowbreak{}escuro.\allowbreak{}svg} /
\texttt{.\allowbreak{}png} & logo horizontal para fundo escuro \\
\texttt{bruma-\allowbreak{}logo-\allowbreak{}claro.\allowbreak{}svg} /
\texttt{.\allowbreak{}png} & logo horizontal para fundo claro
(documentos, pôsteres) \\
\texttt{bruma-\allowbreak{}16.\allowbreak{}png} \ldots{}
\texttt{bruma-\allowbreak{}1024.\allowbreak{}png} & ícone em vários
tamanhos \\
\texttt{bruma.\allowbreak{}ico} & ícone do Windows (16--256 px) \\
\end{longtable}

\textbf{Significado:}

\begin{itemize}
\tightlist
\item
  \textbf{linhas onduladas:} bandas de valência (ciano) e de condução
  (violeta), como camadas de névoa;
\item
  \textbf{linha âmbar tracejada:} o nível de Fermi;
\item
  \textbf{ponto âmbar:} um estado de defeito no gap.
\end{itemize}

\textbf{Cores:}

\begin{longtable}[]{@{}ll@{}}
\toprule\noalign{}
Cor & Código \\
\midrule\noalign{}
\endhead
\bottomrule\noalign{}
\endlastfoot
fundo & \texttt{\#070C0B} \\
ciano & \texttt{\#4FE3C1} \\
âmbar & \texttt{\#F5B544} \\
violeta & \texttt{\#8B7CF6} \\
\end{longtable}

A fonte é a Space Grotesk 600, a mesma do site. Ela já está convertida
em contornos nos SVG.

\textbf{Regerar PNG e ICO depois de editar um SVG:}

\begin{Shaded}
\begin{Highlighting}[]
\ExtensionTok{inkscape}\NormalTok{ bruma{-}mark.svg }\AttributeTok{{-}w}\NormalTok{ 256 }\AttributeTok{{-}h}\NormalTok{ 256 }\AttributeTok{{-}o}\NormalTok{ bruma{-}256.png          }\CommentTok{\# um tamanho}
\ExtensionTok{magick}\NormalTok{ bruma{-}16.png bruma{-}24.png bruma{-}32.png bruma{-}48.png bruma{-}64.png bruma{-}128.png bruma{-}256.png bruma.ico}
\end{Highlighting}
\end{Shaded}

Sem o ImageMagick, use o Pillow:

\begin{Shaded}
\begin{Highlighting}[]
\ImportTok{from}\NormalTok{ PIL }\ImportTok{import}\NormalTok{ Image}
\NormalTok{ims }\OperatorTok{=}\NormalTok{ [Image.}\BuiltInTok{open}\NormalTok{(}\SpecialStringTok{f"bruma{-}}\SpecialCharTok{\{}\NormalTok{s}\SpecialCharTok{\}}\SpecialStringTok{.png"}\NormalTok{) }\ControlFlowTok{for}\NormalTok{ s }\KeywordTok{in}\NormalTok{ (}\DecValTok{16}\NormalTok{, }\DecValTok{24}\NormalTok{, }\DecValTok{32}\NormalTok{, }\DecValTok{48}\NormalTok{, }\DecValTok{64}\NormalTok{, }\DecValTok{128}\NormalTok{, }\DecValTok{256}\NormalTok{)]}
\NormalTok{ims[}\OperatorTok{{-}}\DecValTok{1}\NormalTok{].save(}\StringTok{"bruma.ico"}\NormalTok{, sizes}\OperatorTok{=}\NormalTok{[im.size }\ControlFlowTok{for}\NormalTok{ im }\KeywordTok{in}\NormalTok{ ims])}
\end{Highlighting}
\end{Shaded}

\begin{center}\rule{0.5\linewidth}{0.5pt}\end{center}

\hypertarget{parte-h-manutenuxe7uxe3o-e-novas-versuxf5es}{%
\section{Parte H: manutenção e novas
versões}\label{parte-h-manutenuxe7uxe3o-e-novas-versuxf5es}}

\hypertarget{h1.-lanuxe7ar-uma-versuxe3o-nova-do-bruma-roteiro}{%
\subsection{H1. Lançar uma versão nova do Bruma
(roteiro)}\label{h1.-lanuxe7ar-uma-versuxe3o-nova-do-bruma-roteiro}}

\begin{enumerate}
\def\labelenumi{\arabic{enumi}.}
\item
  Faça as mudanças no código e teste:
  \texttt{python3 bruma.\allowbreak{}py <seed>},
  \texttt{python3 bruma\_\allowbreak{}gui.\allowbreak{}py} e o
  \texttt{index.\allowbreak{}html} no navegador.
\item
  Troque a versão em todos os arquivos de uma vez:

\begin{Shaded}
\begin{Highlighting}[]
\ExtensionTok{python3}\NormalTok{ tools/versao.py            }\CommentTok{\# mostra a versão atual em cada arquivo}
\ExtensionTok{python3}\NormalTok{ tools/versao.py 1.1.0      }\CommentTok{\# aplica a 1.1.0 e a data de hoje}
\end{Highlighting}
\end{Shaded}
\item
  Recompile o tutorial (F1) e copie o PDF para a raiz e para
  \texttt{docs/\allowbreak{}}.
\item
  Envie e publique:

\begin{Shaded}
\begin{Highlighting}[]
\FunctionTok{git}\NormalTok{ commit }\AttributeTok{{-}am} \StringTok{"Versão 1.1.0"}
\FunctionTok{git}\NormalTok{ tag v1.1.0}
\FunctionTok{git}\NormalTok{ push}
\FunctionTok{git}\NormalTok{ push }\AttributeTok{{-}{-}tags}
\end{Highlighting}
\end{Shaded}

  Isso gera o Windows (D1), o PyPI (E1, se configurado) e o DOI (E2, se
  configurado).
\item
  Atualize o site (B3):

\begin{Shaded}
\begin{Highlighting}[]
\FunctionTok{bash}\NormalTok{ tools/copiar\_para\_site.sh ../brunohdmacedo.github.io}
\end{Highlighting}
\end{Shaded}
\end{enumerate}

Use \textbf{versionamento semântico}:

\begin{longtable}[]{@{}ll@{}}
\toprule\noalign{}
Mudança & Exemplo \\
\midrule\noalign{}
\endhead
\bottomrule\noalign{}
\endlastfoot
correção & 1.0.0 → 1.0.1 \\
recurso novo compatível & 1.0.0 → 1.1.0 \\
mudança que quebra o uso antigo & 1.0.0 → 2.0.0 \\
\end{longtable}

\hypertarget{h2.-mudar-os-seus-dados-foto-links-afiliauxe7uxe3o}{%
\subsection{H2. Mudar os seus dados (foto, links,
afiliação)}\label{h2.-mudar-os-seus-dados-foto-links-afiliauxe7uxe3o}}

\begin{longtable}[]{@{}
  >{\raggedright\arraybackslash}p{(\columnwidth - 2\tabcolsep) * \real{0.5000}}
  >{\raggedright\arraybackslash}p{(\columnwidth - 2\tabcolsep) * \real{0.5000}}@{}}
\toprule\noalign{}
\begin{minipage}[b]{\linewidth}\raggedright
Dado
\end{minipage} & \begin{minipage}[b]{\linewidth}\raggedright
Onde mudar
\end{minipage} \\
\midrule\noalign{}
\endhead
\bottomrule\noalign{}
\endlastfoot
Nome, afiliação, links, foto, citação (linha de comando e janela) &
\texttt{bruma.\allowbreak{}py}: constantes \texttt{AUTOR},
\texttt{AFILIACAO}, \texttt{LINKS}, \texttt{FOTO\_\allowbreak{}URL},
\texttt{CITACAO\_\allowbreak{}ABNT},
\texttt{CITACAO\_\allowbreak{}BIBTEX} \\
Mesmos dados na versão web & \texttt{index.\allowbreak{}html} do Bruma:
objetos \texttt{AUTOR} e \texttt{BRUMA} e a constante
\texttt{BIB\_\allowbreak{}BRUMA} (procure por ``Sobre o Bruma'') \\
Texto da janela (afiliação e pesquisa na aba Sobre) &
\texttt{bruma\_\allowbreak{}gui.\allowbreak{}py}: função
\texttt{\_\allowbreak{}build\_\allowbreak{}about} \\
Site pessoal & \texttt{index.\allowbreak{}html} da raiz do site (A3) \\
Citações & \texttt{bruma.\allowbreak{}bib},
\texttt{CITATION.\allowbreak{}cff} \\
\end{longtable}

A \textbf{foto} vem do Lattes (\texttt{FOTO\_\allowbreak{}URL}). Se você
trocar a foto no Lattes, todas as versões passam a mostrar a nova; o
instalador do Windows só a atualiza na próxima compilação.

\hypertarget{h3.-cuxf3pias-de-seguranuxe7a}{%
\subsection{H3. Cópias de
segurança}\label{h3.-cuxf3pias-de-seguranuxe7a}}

O GitHub guarda todo o histórico. Mesmo assim, de tempos em tempos,
baixe um zip de cada repositório (\textbf{Code → Download ZIP}) e guarde
junto com os seus cálculos.

\begin{center}\rule{0.5\linewidth}{0.5pt}\end{center}

\hypertarget{parte-i-soluuxe7uxe3o-de-problemas}{%
\section{Parte I: solução de
problemas}\label{parte-i-soluuxe7uxe3o-de-problemas}}

\begin{longtable}[]{@{}
  >{\raggedright\arraybackslash}p{(\columnwidth - 4\tabcolsep) * \real{0.3333}}
  >{\raggedright\arraybackslash}p{(\columnwidth - 4\tabcolsep) * \real{0.3333}}
  >{\raggedright\arraybackslash}p{(\columnwidth - 4\tabcolsep) * \real{0.3333}}@{}}
\toprule\noalign{}
\begin{minipage}[b]{\linewidth}\raggedright
Problema
\end{minipage} & \begin{minipage}[b]{\linewidth}\raggedright
Causa provável
\end{minipage} & \begin{minipage}[b]{\linewidth}\raggedright
Solução
\end{minipage} \\
\midrule\noalign{}
\endhead
\bottomrule\noalign{}
\endlastfoot
O site não mudou depois do push & publicação em andamento ou cache &
veja a aba \textbf{Actions}; recarregue com Ctrl+F5 \\
``There isn't a GitHub Pages site here'' & Pages desativado &
\textbf{Settings → Pages → Source} (A2) \\
Domínio perdeu o HTTPS & \texttt{CNAME} apagado & restaure o
\texttt{CNAME} com \texttt{brunohdmacedo.\allowbreak{}tech} \\
Citações não aparecem no site & workflow nunca rodou & \textbf{Actions →
Atualizar citações → Run workflow} (A5) \\
O workflow de citações falha no push & token sem permissão de escrita &
\textbf{Settings → Actions → General → Read and write} \\
Só aparecem números do OpenAlex & \texttt{SERPAPI\_\allowbreak{}KEY}
ausente ou cota esgotada & confira o secret; veja \texttt{avisos} em
\texttt{data/\allowbreak{}citacoes.\allowbreak{}json} \\
\texttt{No module named tkinter} & Tk não instalado &
\texttt{sudo apt install python3-\allowbreak{}tk} (Linux) \\
\texttt{bruma:\allowbreak{} command not found} &
\texttt{\textasciitilde{}/\allowbreak{}.\allowbreak{}local/\allowbreak{}bin}
fora do PATH &
\texttt{export PATH="\$HOME/\allowbreak{}.\allowbreak{}local/\allowbreak{}bin:\allowbreak{}\$PATH"} \\
Erro de display no cluster & matplotlib tentando abrir janela &
\texttt{export MPLBACKEND=Agg} (o Bruma já usa Agg) \\
``O Windows protegeu o computador'' & programa sem assinatura digital &
\textbf{Mais informações → Executar assim mesmo} (D4) \\
O workflow do Windows falhou no Inno Setup & instalação do Chocolatey
falhou & rode de novo (\textbf{Re-run jobs}) \\
Foto aparece como ``BM'' & sem internet ou Lattes fora do ar & normal; a
foto aparece quando houver conexão \\
\texttt{xelatex} reclama de \texttt{polyglossia} & TeX incompleto &
instale o TeX Live completo ou use o Overleaf \\
Referência ``??'' no PDF & BibTeX não rodou & rode
\texttt{xelatex → bibtex → xelatex → xelatex} \\
\texttt{pip install bruma-\allowbreak{}dft} não acha o pacote & ainda
não publicado no PyPI & Parte E1, ou use \texttt{bruma.\allowbreak{}py}
direto \\
\end{longtable}

\begin{center}\rule{0.5\linewidth}{0.5pt}\end{center}

\hypertarget{apuxeandice-comandos-ruxe1pidos-e-mapa-de-arquivos}{%
\section{Apêndice: comandos rápidos e mapa de
arquivos}\label{apuxeandice-comandos-ruxe1pidos-e-mapa-de-arquivos}}

\hypertarget{comandos-ruxe1pidos}{%
\subsection{Comandos rápidos}\label{comandos-ruxe1pidos}}

\begin{Shaded}
\begin{Highlighting}[]
\CommentTok{\# Site}
\BuiltInTok{cd}\NormalTok{ brunohdmacedo.github.io }\KeywordTok{\&\&} \FunctionTok{git}\NormalTok{ add }\AttributeTok{{-}A} \KeywordTok{\&\&} \FunctionTok{git}\NormalTok{ commit }\AttributeTok{{-}m} \StringTok{"…"} \KeywordTok{\&\&} \FunctionTok{git}\NormalTok{ push}
\ExtensionTok{python3} \AttributeTok{{-}m}\NormalTok{ http.server 8000                          }\CommentTok{\# testar localmente}

\CommentTok{\# Bruma: análise}
\ExtensionTok{python3}\NormalTok{ bruma.py mgh2\_Fe }\AttributeTok{{-}e}\NormalTok{ Mg,H,Fe                  }\CommentTok{\# cluster}
\ExtensionTok{python3}\NormalTok{ bruma\_gui.py                                 }\CommentTok{\# janela}
\ExtensionTok{python3}\NormalTok{ bruma.py }\AttributeTok{{-}{-}sobre}                             \CommentTok{\# autor e citação}

\CommentTok{\# Bruma: manutenção}
\ExtensionTok{python3}\NormalTok{ tools/versao.py 1.1.0                        }\CommentTok{\# nova versão em todos os arquivos}
\FunctionTok{bash}\NormalTok{ tools/copiar\_para\_site.sh ../brunohdmacedo.github.io}
\FunctionTok{git}\NormalTok{ tag v1.1.0 }\KeywordTok{\&\&} \FunctionTok{git}\NormalTok{ push }\AttributeTok{{-}{-}tags}                    \CommentTok{\# Windows, PyPI e DOI}

\CommentTok{\# LaTeX}
\ExtensionTok{xelatex}\NormalTok{ bruma{-}tutorial }\KeywordTok{\&\&} \ExtensionTok{bibtex}\NormalTok{ bruma{-}tutorial }\KeywordTok{\&\&} \ExtensionTok{xelatex}\NormalTok{ bruma{-}tutorial }\KeywordTok{\&\&} \ExtensionTok{xelatex}\NormalTok{ bruma{-}tutorial}
\end{Highlighting}
\end{Shaded}

\hypertarget{mapa-do-reposituxf3rio-bruma-dft}{%
\subsection{\texorpdfstring{Mapa do repositório
\texttt{bruma-\allowbreak{}dft}}{Mapa do repositório }}\label{mapa-do-reposituxf3rio-bruma-dft}}

\begin{verbatim}
bruma-dft/
├── bruma.py                     linha de comando (motor de análise)
├── bruma_gui.py                 interface gráfica (Tkinter)
├── index.html · tutorial.html   versão web e tutorial
├── bruma-tutorial.tex / .pdf    tutorial em LaTeX/PDF
├── bruma.bib · CITATION.cff     citações
├── pyproject.toml               pacote PyPI (comandos bruma e bruma-gui)
├── requirements.txt
├── assets/                      logo e ícones
├── img/ · docs/img/             figuras
├── docs/MANUAL.md               este manual (GitHub)
├── docs/manual/                 este manual em LaTeX e PDF
├── examples/                    jobs SLURM e PBS
├── tools/versao.py              troca a versão em todos os arquivos
├── tools/copiar_para_site.sh    copia a versão web para o site
├── packaging/windows/           Inno Setup (bruma.iss) e build_windows.bat
└── .github/workflows/           windows.yml (exe) e pypi.yml (PyPI)
\end{verbatim}

\hypertarget{mapa-do-reposituxf3rio-do-site}{%
\subsection{Mapa do repositório do
site}\label{mapa-do-reposituxf3rio-do-site}}

\begin{verbatim}
brunohdmacedo.github.io/
├── index.html                   página principal
├── site/js/main.js              animações e texto digitado
├── site/js/citacoes.js          citações automáticas
├── data/citacoes.json           gerado pelo workflow (não editar à mão)
├── scripts/atualizar_citacoes.py
├── .github/workflows/citacoes.yml
├── bruma/                       cópia da versão web do Bruma
├── castep/                      redirecionamento para /bruma/
├── CITACOES.md · MANUAL.md      guias
├── sitemap.xml · robots.txt · CNAME
└── (páginas e pastas antigas mantidas: v1, v2, v3, portfolio, project-*, 1/, 2/, 3/, robenedet/, cs/…)
\end{verbatim}

\begin{center}\rule{0.5\linewidth}{0.5pt}\end{center}

\emph{Bruma 1.0.0 · Bruno Henrique Dourado Macedo ·
\url{https://brunohdmacedo.tech/bruma/} · ORCID 0000-0002-8152-0950}

\end{document}
